\documentclass[11pt]{article}
\usepackage[a4paper,margin=24mm]{geometry}
\usepackage[T1]{fontenc}
\usepackage{amsmath,amssymb,amsthm,mathtools}
\usepackage[hidelinks,hypertexnames=false]{hyperref}
\setlength{\emergencystretch}{2em}
\newcommand{\C}{\mathbb C}\newcommand{\R}{\mathbb R}
\newcommand{\tr}{\operatorname{tr}}\newcommand{\Res}{\operatorname{Res}}
\newcommand{\id}{\operatorname{Id}}\newcommand{\dd}{\dagger_G}
\newtheorem{theorem}{Theorem}[section]
\newtheorem{lemma}[theorem]{Lemma}
\title{Native spectral action along the outgoing arithmetic relation\\
Exact cyclic coefficients, full determinants and metric curvature}
\author{}\date{20 September 2026}
\begin{document}\maketitle

\section{Objects, source use and scope}

Connes and van Suijlekom \cite{CV}, Section \texttt{sect:sa},
apply divided differences to a trace action with a rank-one,
not necessarily selfadjoint perturbation. Here that calculation
is carried out for the original attained arithmetic quotient
and its outgoing relation. Every trace below is finite.
The test function $F$ is an arbitrary polynomial, so every
expansion in its perturbation parameter terminates. No
infinite trace, spectral limit or Riemann-hypothesis endpoint
is inferred.

Retain the actual packet, with complete zero orders, and
the original arithmetic source:
\begin{equation}\tag{SA1}\label{SA1}
\begin{gathered}
h(s)=\prod_{\rho\in Z}(s-\rho)^{m_\rho},\qquad
g(s)=2\xi(s)=s(s-1)\pi^{-s/2}\Gamma(s/2)\zeta(s),\\
\mathcal MF_h=g/h,\qquad
w_h(v)=\frac{|(g/h)(1/2+iv)|^2}{2\pi},\qquad
m(u)=w_h^{*k}(u),\qquad c=k/2,\quad S=c+iu.
\end{gathered}
\end{equation}
The original nonempty packet is stable under conjugation
and $s\mapsto1-s$. The measure is the one in \cite{Native},
NS1--NS3, not a selected comparison measure. In particular
it is even, positive almost everywhere, has every polynomial
moment, and has mass
$\mathfrak m=(\int_\R w_h)^k$, which is retained.
Let $Q_n$ be its real monic orthogonal polynomials and put
\begin{equation}\tag{SA2}\label{SA2}
\omega_n=\int_\R Q_n(u)^2m(u)\,du,\quad
p_n(S)=i^nQ_n((S-c)/i),\quad
\omega_0=\mathfrak m,\quad
a_n=\sqrt{\omega_{n+1}/\omega_n}.
\end{equation}
All $\omega_n$ are strictly positive: a nonzero polynomial
cannot vanish almost everywhere for this measure.
Gram--Schmidt and evenness give
$Q_n(-u)=(-1)^nQ_n(u)$ and
$uQ_n=Q_{n+1}+(\omega_n/\omega_{n-1})Q_{n-1}$.
For the last assertion, pair with $Q_j$ for $j<n-1$,
transfer $u$ to $Q_j$, and use orthogonality; parity removes
the $Q_n$ coefficient and the $Q_{n-1}$ pairing gives
the displayed coefficient.

Let $C=\C[S]/(\chi)$ be the original cyclic algebra,
$q=\deg\chi$, and $M_S$ multiplication by $S$.
The full annihilator, including sum collisions, is
\begin{equation}\tag{SA3}\label{SA3}
\chi(S)=\prod_\kappa(S-\kappa)^{e_\kappa},
\qquad
e_\kappa=\max_{\sum_j\rho_j=\kappa}
\left(1+\sum_{j=1}^k(m_{\rho_j}-1)\right).
\end{equation}
Indeed the local tensor ring has coordinates $z_j$ with
$z_j^{m_{\rho_j}}=0$, and the highest nonzero power of
$\sum z_j$ has coefficient
$(\sum(m_{\rho_j}-1))!/\prod(m_{\rho_j}-1)!$ on
$\prod z_j^{m_{\rho_j}-1}$. This proves the exponent and,
by taking the maximum at equal sums, the full cyclic kernel.
Set
\begin{equation}\tag{SA4}\label{SA4}
M=(M_S-cI)/i,\quad
\chi_u(z)=i^{-q}\chi(c+iz)=\det(zI-M),\quad b_n=[p_n]_\chi.
\end{equation}
The determinant assertion follows from the companion
matrix of multiplication in the increasing-power basis.
Fix $L\geq q-1$ and use this original remainder basis
throughout unless a coordinate map is explicitly stated.
The attained metric and lifting maps are
\begin{equation}\tag{SA5}\label{SA5}
E=\left[\frac{i^{-n}b_n}{\sqrt{\omega_n}}\right]_{n=0}^L,\quad
G=(EE^*)^{-1},\quad R=E^*G,\quad ER=I,\quad R^*R=G.
\end{equation}
The first $q$ monic columns make $E$ onto. Every lift of
$v$ is $Rv+z$, $Ez=0$, and
$\|Rv+z\|^2=v^*Gv+\|z\|^2$, since
$(Rv)^*z=v^*GEz=0$. Thus $G$ is precisely the least-norm
quotient metric, including the original mass.

Let $J_L$ be the $(L+1)$-square Jacobi matrix in the
orthonormal coordinates $Q_n/\sqrt{\omega_n}$.
Define the following \emph{different} operators on the
\emph{same} algebra and prove their exact relation:
\begin{equation}\tag{SA6}\label{SA6}
A=EJ_LR=M+Q,\qquad
Q=r\ell,\qquad
r=\frac{i}{\omega_L}b_{L+1},\quad \ell=b_L^*G.
\end{equation}
To check the sign, multiplication by $u$ has its extra
outgoing column $a_Le_{L+1}e_L^*$. Removing that column
after reduction and lifting subtracts
\[
a_L\frac{i^{-(L+1)}b_{L+1}}{\sqrt{\omega_{L+1}}}
\frac{i^Lb_L^*G}{\sqrt{\omega_L}}
=-\frac{i}{\omega_L}b_{L+1}b_L^*G.
\]
This proves \eqref{SA6}. Moreover $GA=R^*J_LR$ is
Hermitian, so $A^{\dd}=A$, where $X^{\dd}=G^{-1}X^*G$.
Conjugation by $G^{1/2}$ proves that $A$ has real spectrum
and a complete $G$-orthonormal eigenbasis.

The parity map $\Gamma[P(S)]=[P(k-S)]$ is well-defined
on $C$ by packet symmetry. Its entries are
$\Gamma_{rj}=\binom jr k^{j-r}(-1)^r$, $0\leq r\leq j<q$.
It satisfies
\begin{equation}\tag{SA7}\label{SA7}
\Gamma^2=I,\quad \Gamma^*G\Gamma=G,\quad
\Gamma A\Gamma=-A,\quad \Gamma Q\Gamma=-Q,\quad
\ell r=0,\quad Q^2=0.
\end{equation}
For proof, $\Gamma b_n=(-1)^nb_n$, and hence
$\Gamma E=E\operatorname{diag}((-1)^n)$.
Use \eqref{SA5} and the odd Jacobi matrix to prove the
first four identities. Opposite parity gives
$b_L^*Gb_{L+1}=-b_L^*Gb_{L+1}=0$, proving the last two.
This is a calculation of the pairing, not removal of a
cross term.

\section{Polynomial first principles and every repeated node}

For a polynomial $f$ define divided differences at any
nodes, including repetitions, by linearity and
\begin{equation}\tag{SA8}\label{SA8}
(x^m)[x_1,\ldots,x_n]=
\begin{cases}
\displaystyle\sum_{\substack{a_1,\ldots,a_n\geq0\\
 \sum a_j=m-n+1}}\prod_{j=1}^n x_j^{a_j},&m\geq n-1,\\
0,&m<n-1.
\end{cases}
\end{equation}
This agrees with the recursive divided difference:
subtract the complete homogeneous sums with first or last
variable omitted and factor the difference of the two
powers by the difference of those variables. Induction
on $n$ starts with $f[x_1]=f(x_1)$.
The resulting polynomial identity also proves the unique
continuous extension to coincident nodes. In particular
$f[x,\ldots,x]=f^{(n-1)}(x)/(n-1)!$.
The elementary Laurent identity behind this definition is
\begin{equation}\tag{SA9}\label{SA9}
\Res_\infty^+\frac{f(z)}
 {\prod_{j=1}^n(z-x_j)}
=f[x_1,\ldots,x_n],
\end{equation}
where $\Res_\infty^+$ denotes the coefficient of $z^{-1}$
in the expansion at infinity, not its negative.
Expand each $(z-x_j)^{-1}$ as a geometric series to obtain
\eqref{SA8}, proving \eqref{SA9} also at repeated nodes.

\begin{theorem}[Exact cyclic trace coefficient]\label{thm:cyclic}
Let $A$ be the operator in \eqref{SA6}, let $V$ be any
operator on the same finite space, and let
$A=\sum_{\lambda\in\Lambda}\lambda P_\lambda$ be its
decomposition into distinct real eigenvalues, with the
complete spectral projections. For every polynomial $F$
and $n\geq1$,
\begin{equation}\tag{SA10}\label{SA10}
\left.\frac{d^n}{dt^n}\tr F(A+tV)\right|_{t=0}
=(n-1)!\sum_{\lambda_1,\ldots,\lambda_n\in\Lambda}
F'[\lambda_1,\ldots,\lambda_n]\,
\tr(P_{\lambda_1}V\cdots P_{\lambda_n}V).
\end{equation}
Every multiplicity of $A$ is retained inside its projector.
\end{theorem}
\begin{proof}
First take $F(x)=x^m$. The product rule and cyclicity give
$d\,\tr(A+tV)^m/dt=m\,\tr((A+tV)^{m-1}V)$:
the $m$ differentiated words have equal trace by cyclic
permutation. The coefficient of $t^{n-1}$ in the right
side is
\[
m\sum_{a_1+\cdots+a_n=m-n}
\tr(A^{a_1}V A^{a_2}V\cdots A^{a_n}V).
\]
Insert the complete projector decomposition of each power.
The sum over exponents is exactly the divided difference
of $F'=mx^{m-1}$ on $n$ nodes, by \eqref{SA8}.
Since differentiating a coefficient of $t^n$ multiplies
it by $n$, the coefficient in the trace is $1/n$ times
this sum. Multiplication by $n!$ proves \eqref{SA10}.
Linearity proves the result for all $F$. Every operation
was a finite polynomial operation, so repetitions and
zero eigenvalues require no limiting argument.
\end{proof}

The statement of \cite{CV}, Lemma \texttt{lem:sa}, displays
$n!$ rather than $(n-1)!$ for the sum in \eqref{SA10}.
The proof there instead ends with the correct repeated-node
expression
$n!\sum\tr(P_1V\cdots P_nV)F[\lambda_1,\ldots,\lambda_n,\lambda_1]$.
Here is the exact passage between them. By \eqref{SA9},
differentiation of $\prod(z-\lambda_j)^{-1}$ and extraction
of the coefficient of a derivative of a Laurent series give
\begin{equation}\tag{SA11}\label{SA11}
F'[\lambda_1,\ldots,\lambda_n]
=\sum_{j=1}^n
F[\lambda_1,\ldots,\lambda_n,\lambda_j].
\end{equation}
In the full cyclic sum the $n$ summands on the right
have equal summed values, by cyclic relabelling.
Thus the repeated-node expression is $n!/n=(n-1)!$
times the $F'$ expression. The source's subsequent
second-derivative formula and Proposition
\texttt{prop:der-sa} use the latter, correct second
derivative. For an immediate coefficient test, take
dimension one, $A=0$, $V=1$, $F(x)=x^2$: the second
derivative is $2$, whereas the displayed $n!F'[0,0]$
would give $4$. No other assertion of that source is
needed for this finite correction.

\section{The complete native path, determinant and resolvent}

Return to the actual $Q=r\ell$ in \eqref{SA6} and define
\begin{equation}\tag{SA12}\label{SA12}
T(t)=A-tQ,\quad T(0)=A,\quad T(1)=M,\qquad
d(z)=\det(zI-A),\quad
\kappa(z)=\ell(zI-A)^{-1}r.
\end{equation}
The parameter $t$ is an algebraic perturbation parameter;
all identities of polynomials below hold for every complex
$t$. The $G$-adjoint statements later use real $t$.
Put $n(z)=\ell\,\operatorname{adj}(zI-A)r$.
The full, uncancelled characteristic polynomial is
\begin{equation}\tag{SA13}\label{SA13}
\begin{split}
p_t(z)&=\det(zI-T(t))=d(z)+t n(z)
       =d(z)(1+t\kappa(z))\\
      &=(1-t)d(z)+t\chi_u(z),\qquad
n(z)=\chi_u(z)-d(z).
\end{split}
\end{equation}
For proof, away from the spectrum of $A$, factor
$zI-T=(zI-A)(I+t(zI-A)^{-1}r\ell)$ and use
$\det(I+v\ell)=1+\ell v$. A basis beginning with $v$
proves this last identity when $v\ne0$; the zero case is
immediate. Multiplication by $d$ gives a polynomial
identity everywhere. Setting $t=1$ gives the last line.
No common factor or invisible eigenvalue was removed.

For all points where the displayed inverses exist, direct
multiplication gives
\begin{equation}\tag{SA14}\label{SA14}
(zI-T(t))^{-1}
=(zI-A)^{-1}
-\frac{t(zI-A)^{-1}r\ell(zI-A)^{-1}}
       {1+t\kappa(z)}.
\end{equation}
It is also an identity of rational matrices.
Taking its trace and using
$\kappa'(z)=-\ell(zI-A)^{-2}r$ yields
\begin{equation}\tag{SA15}\label{SA15}
\tr(zI-T(t))^{-1}-\tr(zI-A)^{-1}
=\frac{t\kappa'(z)}{1+t\kappa(z)}
=\partial_z\log(1+t\kappa(z)).
\end{equation}
The logarithm is unambiguously its expansion tending to
zero at infinity. Its derivative is rational, so the
identity has no branch dependence.

Write $w_\lambda=\ell P_\lambda r$. The full scalar
resolvent and its moments are
\begin{equation}\tag{SA16}\label{SA16}
\kappa(z)=\sum_{\lambda\in\Lambda}
\frac{w_\lambda}{z-\lambda}
=\sum_{j\geq0}\frac{\ell A^jr}{z^{j+1}},
\qquad \sum_\lambda w_\lambda=\ell r=0.
\end{equation}
Spectral projectors are in the original metric. Explicitly,
if $U_\lambda$ contains a basis of the eigenspace, then
$P_\lambda=U_\lambda(U_\lambda^*GU_\lambda)^{-1}
U_\lambda^*G$. Distinct eigenspaces are $G$-orthogonal
by selfadjointness, proving this expression and the
completeness used above.

\begin{theorem}[Trace, determinant and residue are the same calculation]
For every polynomial $F$ of degree $m$,
\begin{equation}\tag{SA17}\label{SA17}
\begin{split}
\tr F(T(t))-\tr F(A)
&=-\Res_\infty^+ F'(z)\log(1+t\kappa(z))\\
&=\sum_{n=1}^{m}\frac{(-t)^n}{n}
\Res_\infty^+ F'(z)\kappa(z)^n\\
&=\sum_{n=1}^{m}\frac{(-t)^n}{n}
\sum_{\lambda_1,\ldots,\lambda_n}
F'[\lambda_1,\ldots,\lambda_n]\prod_{j=1}^n w_{\lambda_j}.
\end{split}
\end{equation}
\end{theorem}
\begin{proof}
Expand the resolvent at infinity as
$\sum_{j\geq0}T^jz^{-j-1}$. The coefficient of $z^{-1}$
in $F(z)(zI-T)^{-1}$ is $F(T)$.
Use \eqref{SA15} and take its trace to obtain the first
line after integration by parts in Laurent series:
the coefficient of $z^{-1}$ of any derivative is zero.
Since $\kappa=O(z^{-1})$, terms $n>m$ cannot contribute.
The logarithm series therefore gives the second line
as a finite polynomial identity in $t$, even though its
auxiliary series is local at infinity. Insert
\eqref{SA16} and apply \eqref{SA9} for the third line.
Alternatively each trace cycle in \eqref{SA10} for
$V=-r\ell$ is $(-1)^n\prod_j\ell P_{\lambda_j}r$.
Thus the two independent finite derivations agree.
\end{proof}

The original parity strengthens these formulas:
\begin{equation}\tag{SA18}\label{SA18}
w_{-\lambda}=-w_\lambda,\quad w_0=0,\quad
\kappa(-z)=\kappa(z),\quad
\kappa(z)=\sum_{j\geq0}\frac{\ell A^{2j+1}r}{z^{2j+2}},
\quad \deg n\leq q-2.
\end{equation}
Indeed $P_{-\lambda}=\Gamma P_\lambda\Gamma$,
$\ell\Gamma=(-1)^L\ell$, and
$\Gamma r=(-1)^{L+1}r$, proving the weights.
The Laurent statements follow. In particular
\begin{equation}\tag{SA19}\label{SA19}
\deg_t\tr F(T(t))\leq\lfloor m/2\rfloor,\qquad
\tr T(t)^{2j+1}=0.
\end{equation}
For the degree bound use $\kappa=O(z^{-2})$ in
\eqref{SA17}; a $z^{-1}$ term needs $m\geq2n$.
For odd traces, $\Gamma T(t)\Gamma=-T(t)$ and cyclic
invariance of trace prove the vanishing. Nilpotence
has not made $T(t)$ isospectral.

\section{First and second derivatives, with their exact cancellations}

Let $\Phi_F(t)=\tr F(T(t))$ and put $\gamma_j=\ell A^jr$.
The first two derivatives are
\begin{equation}\tag{SA20}\label{SA20}
\Phi_F'(0)=-\ell F'(A)r,\qquad
\Phi_F''(0)=
\sum_{\lambda,\mu}w_\lambda w_\mu F'[\lambda,\mu].
\end{equation}
The second expression has no conjugates on its weights:
the perturbation is not $G$-selfadjoint. It is not a sum
of positive squared matrix entries.
The full native coordinate change
\begin{equation}\tag{SA21}\label{SA21}
\begin{gathered}
T_{rj}^{\,\mathrm{coord}}=\binom jr c^{j-r}i^r,\quad
G_u=((T^{\mathrm{coord}})^{-1})^*G
                 (T^{\mathrm{coord}})^{-1},\\
d_n=[Q_n]_{\chi_u}=i^{-n}T^{\mathrm{coord}}b_n,\quad
Q_u=-d_{L+1}d_L^*G_u/\omega_L,\quad
A_u=M_u+Q_u
\end{gathered}
\end{equation}
has a real positive $G_u$ and real $A_u,Q_u$.
The triangular coordinate formula expands $(c+iu)^j$;
inserting it into \eqref{SA5} and \eqref{SA6} proves
every displayed sign. Thus the scalars $\gamma_j$ and
$w_\lambda$ are real for the actual symmetric packet.
This conclusion follows from the map, not from discarding
the original complex phases.

For any scalar $a$, the exact polynomial identity
\begin{equation}\tag{SA22}\label{SA22}
F'[x,y]-F'[x,a]-F'[a,y]+F''(a)
=(x-a)(y-a)F'[x,y,a,a]
\end{equation}
holds at all nodes, including coincidences.
Subtract the fractions in \eqref{SA9}:
\[
\frac1{(z-x)(z-y)}
-\frac1{(z-x)(z-a)}
-\frac1{(z-a)(z-y)}
+\frac1{(z-a)^2}
=\frac{(x-a)(y-a)}
       {(z-x)(z-y)(z-a)^2}.
\]
Multiplication by $F'$ and extraction proves
\eqref{SA22}. Since $\sum w_\lambda=0$, it gives
\begin{equation}\tag{SA23}\label{SA23}
\Phi_F''(0)=
\sum_{\lambda,\mu}w_\lambda w_\mu
(\lambda-a)(\mu-a)F'[\lambda,\mu,a,a].
\end{equation}
Consequently the second derivative vanishes for every
polynomial of degree at most three. Two explicit full
trace formulas are
\begin{equation}\tag{SA24}\label{SA24}
\begin{split}
\tr T(t)^2&=\tr A^2-2t\gamma_1,\\
\tr T(t)^4&=\tr A^4-4t\gamma_3+2t^2\gamma_1^2 .
\end{split}
\end{equation}
These follow immediately from \eqref{SA17}, or from
expanding the words and using $Q^2=0$. They retain the
mixed products $AQ$, which do not vanish with $Q^2$.
Neither the equality of first and second derivatives
in \cite{CV}, Proposition \texttt{prop:der-sa}, nor the
specific relation $QD\xi=-\beta$ used in its proof has
been substituted for \eqref{SA20}. Here the exact
receiving function is the calculated $\kappa$.

There is a finite quantitative connection to the existing
control even for every higher trace derivative. Set
\begin{equation}\tag{SA25}\label{SA25}
\epsilon=\|Q\|_G
=\frac{\sqrt{(b_L^*Gb_L)(b_{L+1}^*Gb_{L+1})}}
 {\omega_L},\qquad
\sum_\lambda|w_\lambda|\leq\epsilon .
\end{equation}
For proof choose a $G$-orthonormal eigenbasis.
Write $r_j$ for the coordinates of $r$ and $\ell_j$ for
those of its row. Triangle inequality within repeated
eigenspaces and Cauchy--Schwarz give
$\sum_\lambda|w_\lambda|
\leq(\sum|r_j|^2)^{1/2}(\sum|\ell_j|^2)^{1/2}$.
An outer product has norm the product of these norms,
which is the first expression in \eqref{SA25}.
If $I$ is the smallest closed real interval containing
$\operatorname{spec}A$, then for real or complex $F$,
\begin{equation}\tag{SA26}\label{SA26}
|\Phi_F^{(n)}(0)|
\leq\epsilon^n\max_{x\in I}|F^{(n)}(x)|.
\end{equation}
The repeated fundamental theorem of calculus gives
$F'[x_1,\ldots,x_n]
=\int_{\Delta_{n-1}}F^{(n)}(\sum s_jx_j)\,ds$.
Its volume is $1/(n-1)!$ (successive integration proves
this by induction). Combine this bound with
\eqref{SA10} and \eqref{SA25}. This is a bound in terms
of the actual $\epsilon$, not an estimate making that
quantity decay.

\section{The exact metric action which recovers epsilon}

For real $t$ define the same-metric squared norm and
Hermitian defect
\begin{equation}\tag{SA27}\label{SA27}
\mathcal H(t)=\tr(T(t)^{\dd}T(t)),\qquad
\mathcal C(t)=i(T(t)-T(t)^{\dd})
=it(Q^{\dd}-Q).
\end{equation}
No choice of new metric occurs. Direct multiplication,
cyclic trace and $A^{\dd}=A$ give
\begin{equation}\tag{SA28}\label{SA28}
\begin{gathered}
\mathcal H(t)=\tr A^2-2t\Re\gamma_1+t^2\epsilon^2,\\
\mathcal H''(t)=2\epsilon^2,\qquad
\tr T(t)^2=\tr A^2-2t\gamma_1,\\
\mathcal H(t)-\Re\tr T(t)^2=t^2\epsilon^2,\qquad
\tr\mathcal C(t)^2=2t^2\epsilon^2.
\end{gathered}
\end{equation}
To prove $\tr(Q^{\dd}Q)=\epsilon^2$, conjugate by
$G^{1/2}$ to the outer product
$xy^*$, where $x=G^{1/2}r$ and
$y=G^{-1/2}\ell^*$. Its squared Hilbert--Schmidt norm is
$\|x\|^2\|y\|^2$, also its squared operator norm.
Moreover $y^*x=\ell r=0$.
When both are nonzero the restriction of
$i(Q^{\dd}-Q)$ in the orthonormal pair
$x/\|x\|,y/\|y\|$, after this exact metric conjugation,
is
$\left(\begin{smallmatrix}0&-i\epsilon\\
i\epsilon&0\end{smallmatrix}\right)$.
It vanishes on their orthogonal complement. This proves
the last equality and the eigenvalues
$+\epsilon,-\epsilon$, with all other eigenvalues zero.
If either vector is zero all displayed quantities are
zero, which proves the degenerate case too.

The two-variable polynomial action gives an equally exact
mixed derivative without taking any complex conjugate of
the independent parameters:
\begin{equation}\tag{SA29}\label{SA29}
\mathcal J(z,w)=\tr((A-wQ^{\dd})(A-zQ))
=\tr A^2-z\gamma_1-w\overline{\gamma_1}+zw\epsilon^2,
\qquad \partial_z\partial_w\mathcal J=\epsilon^2 .
\end{equation}
Thus the vanishing holomorphic second derivative of the
quadratic trace in \eqref{SA24} and the nonzero metric
second derivative in \eqref{SA28} are connected by a
displayed polynomial map. They are not unrelated
observations. The mixed coefficient is exactly the
original arithmetic allowance squared.

One can also express this as a selfadjoint trace action
on a specified doubled space. On $C\oplus C$ with metric
$G\oplus G$, set
\begin{equation}\tag{SA30}\label{SA30}
\mathscr D(t)=
\begin{pmatrix}0&T(t)\\T(t)^{\dd}&0\end{pmatrix},
\qquad
\tr\mathscr D(t)^2=2\mathcal H(t),\qquad
\frac{d^2}{dt^2}\tr\mathscr D(t)^2=4\epsilon^2 .
\end{equation}
The block adjoint proves selfadjointness and block
multiplication proves the trace identities. Its complete
characteristic polynomial, including zero factors, is
\begin{equation}\tag{SA31}\label{SA31}
\det(zI_{2q}-\mathscr D(t))
=\det(z^2I_q-T(t)^{\dd}T(t)).
\end{equation}
For $z\ne0$ this is the Schur determinant with the
$zI_q$ upper-left block; the powers of $z$ cancel
exactly. Polynomial continuation proves it also at
$z=0$. This dilation observes singular values in the
retained metric; its determinant is explicitly
\eqref{SA31}, not the polynomial $p_t$.

\section{Every original S-phase and the physical image}

The original scaling path is
\begin{equation}\tag{SA32}\label{SA32}
B(t)=cI+iT(t)=B(0)-itQ,\qquad
B(0)=M_S-\frac{b_{L+1}b_L^*G}{\omega_L},\quad
B(1)=M_S.
\end{equation}
Its complete determinant and resolvent are
\begin{equation}\tag{SA33}\label{SA33}
\begin{gathered}
\det(sI-B(t))=i^qp_t((s-c)/i)
=(1-t)\det(sI-B(0))+t\chi(s),\\
(sI-B(t))^{-1}
=i^{-1}\big(((s-c)/i)I-T(t)\big)^{-1}.
\end{gathered}
\end{equation}
Factor $sI-B(t)=i(((s-c)/i)I-T(t))$ to verify every
factor. For a polynomial $f$ in the original $S$
coordinate, take $F(u)=f(c+iu)$ in \eqref{SA17}.
Its derivatives satisfy $F^{(n)}(u)=i^nf^{(n)}(c+iu)$,
and its cyclic divided differences satisfy
\begin{equation}\tag{SA34}\label{SA34}
F'[\lambda_1,\ldots,\lambda_n]
=i^n f'[c+i\lambda_1,\ldots,c+i\lambda_n].
\end{equation}
The recursive definition or \eqref{SA8} proves the
factor $i^{n-1}$ from change of nodes and the further
$i$ from differentiating $F$. Thus
\begin{equation}\tag{SA35}\label{SA35}
\left.\frac{d^n}{dt^n}\tr f(B(t))\right|_{t=0}
=(-i)^n(n-1)!
\sum_{\lambda_1,\ldots,\lambda_n}
f'[c+i\lambda_1,\ldots,c+i\lambda_n]\prod_jw_{\lambda_j}.
\end{equation}
The original centered defect is exactly
\begin{equation}\tag{SA36}\label{SA36}
B(t)+B(t)^{\dd}-kI
=\mathcal C(t),\qquad
\operatorname{spec}\mathcal C(t)
=\{t\epsilon,-t\epsilon,0,\ldots,0\}.
\end{equation}
At $t=1$, pairing against an eigenvector of $M_S$
gives $|2\Re\kappa-k|\leq\epsilon$, including a root
whose algebraic multiplicity exceeds one. An
eigenvector exists for each distinct root; its
Rayleigh quotient is $2\Re\kappa-k$.

The complete physical injection is retained as well.
With $E_h=\C[s]/(h)$ and $\upsilon=j_h(g/h)$, define
\begin{equation}\tag{SA37}\label{SA37}
\eta:C\longrightarrow E_h^{\otimes k},\qquad
\eta[P]=\upsilon^{\otimes k}P(s_1+\cdots+s_k).
\end{equation}
Substitution has kernel exactly \eqref{SA3}.
At a zero $\rho$, the constant local coefficient of
$g/h$ equals
$g^{(m_\rho)}(\rho)/
(m_\rho!\prod_{\sigma\ne\rho}(\rho-\sigma)^{m_\sigma})$,
which is nonzero. In a truncated polynomial ring a
series with nonzero constant term is invertible: its
inverse is the finite geometric series after removing
that constant. Consequently $\upsilon$ is a unit in
every local factor, and $\eta$ is injective.
Give precisely $\eta(C)$ the metric
$\langle\eta v,\eta w\rangle=v^*Gw$.
Then $\eta T(t)\eta^{-1}$, $\eta B(t)\eta^{-1}$ and
$\eta Q\eta^{-1}$ have all the traces, determinants
and adjoints proved above. This follows from actual
conjugacy and the displayed transported metric; the
inverse is only on $\eta(C)$.
The physical attained lift of $v$ is
$\mathcal A\sum_{n=0}^L(Rv)_nQ_n/\sqrt{\omega_n}$,
where $\mathcal A P=P(D_1+\cdots+D_k)F_h^{\otimes k}$
and $Q_n$ is evaluated in $u=(S-c)/i$.
Its jet is $\eta v$ by $ER=I$ and its norm squared is
$v^*Gv$ by Parseval and \eqref{SA5}.
Hence the trace path stays connected to the original
arithmetic classes and their actual source lifts.

\section{Exact finite examples and what they establish}

The following examples test finite algebra; neither
is presented as a packet of arithmetic zeros.
For the separate measure $(3/2)1_{[-1,1]}(u)\,du$,
mass $3$, and test quotient $u^2+1$ at $L=1$, one has
\begin{equation}\tag{SA38}\label{SA38}
G=\begin{pmatrix}3&0\\0&1\end{pmatrix},\quad
A=\begin{pmatrix}0&1/3\\1&0\end{pmatrix},\quad
Q=\begin{pmatrix}0&4/3\\0&0\end{pmatrix},\quad
\epsilon^2=16/3,\quad\gamma_1=4/3 .
\end{equation}
Direct multiplication gives $A^{\dd}=A$, $Q^2=0$ and
\begin{equation}\tag{SA39}\label{SA39}
\begin{gathered}
p_t(z)=z^2-\frac13+\frac43t,\qquad
\kappa(z)=\frac{4/3}{z^2-1/3},\\
\tr T(t)^2=\frac23-\frac83t,\qquad
\tr T(t)^4=\frac29-\frac{16}{9}t+\frac{32}{9}t^2,\\
\mathcal H(t)=\frac23-\frac83t+\frac{16}{3}t^2.
\end{gathered}
\end{equation}
This explicitly shows a vanishing quadratic holomorphic
second derivative together with the retained nonzero
metric curvature and moving eigenvalues.

Even every trace polynomial along the path does not by
itself determine $\epsilon$ for arbitrary parity-respecting
finite data with a fixed metric. Here is an exact example
and the missing map, rather than an unsupported
nonidentification. Put $G=I_4$,
$A=\operatorname{diag}(-2,-1,1,2)$,
$\Gamma$ equal to reversal of the four coordinates,
$r=(a,b,b,a)^{\mathsf T}$,
$\ell=(a^{-1},b^{-1},-b^{-1},-a^{-1})$ for positive
$a,b$. Then $\Gamma r=r$, $\ell\Gamma=-\ell$,
$Q^2=0$, and the weights are always $(1,1,-1,-1)$.
Therefore $\kappa$, $p_t$ and every polynomial trace
in \eqref{SA17} are independent of $a,b$. Nevertheless
\begin{equation}\tag{SA40}\label{SA40}
\epsilon^2
=4(a^2+b^2)(a^{-2}+b^{-2})
=\begin{cases}16,&(a,b)=(1,1),\\25,&(a,b)=(2,1).\end{cases}
\end{equation}
The exact map between them is conjugation by the
diagonal matrix commuting with $A$ and $\Gamma$ which
rescales the first and last coordinates. It sends $r$
and $\ell$ contragrediently, retaining $w_\lambda$.
It is not an isometry of the retained $G=I_4$.
Equation \eqref{SA29} measures exactly its change:
the mixed coefficient changes from $16$ to $25$.
This finite example does not assert that both choices
arise from the arithmetic measure \eqref{SA1}.

The same formulas calculate every derivative at the
original endpoint $t=1$, even when $M$ has Jordan blocks.
Indeed \eqref{SA14} gives the exact scalar map
\begin{equation}\tag{SA41}\label{SA41}
\kappa_t(z)=\ell(zI-T(t))^{-1}r
=\frac{\kappa(z)}{1+t\kappa(z)},\qquad
\Phi_F^{(n)}(t)=(-1)^n(n-1)!
\Res_\infty^+F'(z)\kappa_t(z)^n .
\end{equation}
To prove the derivative formula, differentiate the first
line of \eqref{SA17} $n$ times. Differentiation of
$-\log(1+t\kappa)$ gives
$(-1)^n(n-1)!\kappa^n/(1+t\kappa)^n$.
Coefficient extraction at infinity commutes with these
finite polynomial coefficients. This proves the formula
for every complex $t$; it does not require $T(t)$ to be
diagonalizable. Thus higher-order original poles are
retained by the actual rational expression at $t=1$.

Finally the same metric action controls the full
aggregate exterior defect, not only one eigenvalue.
For real $t$, count the eigenvalues $\lambda$ of $B(t)$
with algebraic multiplicity and set
\begin{equation}\tag{SA42}\label{SA42}
\begin{split}
\mathcal L(t)&=\sum_{\Re\lambda>c}
                   (2\Re\lambda-2c),\\
\mathcal L(t)&\leq |t|\epsilon,\qquad
\mathcal L(t)^2\leq
t^2\epsilon^2
=\frac{t^2}{2}\mathcal H''(t)
=\frac{t^2}{4}\frac{d^2}{dt^2}\tr\mathscr D(t)^2.
\end{split}
\end{equation}
Here is the complete finite proof. Take the sum $W$ of
the generalized eigenspaces with $\Re\lambda>c$ and
let $P_W$ be its $G$-orthogonal projection. The space
is invariant, so in a $G$-orthonormal basis starting
with $W$, the upper-left block of $B(t)$ is its
restriction to $W$. Its trace is the sum of the
selected eigenvalues with their full algebraic
multiplicity (put that finite restriction into Jordan
form). Therefore
$\mathcal L(t)=\tr(P_W\mathcal C(t))$.
In a $G$-orthonormal eigenbasis of the Hermitian
$\mathcal C(t)$ each diagonal entry of $P_W$ is in
$[0,1]$. The sum is at most the sum of the positive
eigenvalues of $\mathcal C(t)$, namely $|t|\epsilon$
by \eqref{SA36}. Squaring its nonnegative sides and
using \eqref{SA28}--\eqref{SA30} proves every equality
and inequality in \eqref{SA42}. At $t=1$ this is the
existing aggregate arithmetic spectral defect
bounded by the exact metric curvature of the
calculated outgoing-relation path.

For the exact mass-three example \eqref{SA38}, take
any retained real $c$ and $t\geq0$. Its two $u$-roots
are $\pm\sqrt{(1-4t)/3}$. The original $S$-eigenvalues
are $c$ plus $i$ times those roots, so the aggregate is
\begin{equation}\tag{SA43}\label{SA43}
\mathcal L(t)=
\begin{cases}
0,&0\leq t\leq1/4,\\
2\sqrt{(4t-1)/3},&t>1/4,
\end{cases}
\qquad |t|\epsilon=\frac{4t}{\sqrt3}.
\end{equation}
At $t=1/4$ the repeated eigenvalue is still counted with
its full order and contributes zero. Above it, exactly
one eigenvalue has real part greater than $c$, yielding
the displayed value. Squaring both nonnegative sides
of the bound in \eqref{SA42} reduces here to
$16t-4\leq16t^2$, equivalently
$4(2t-1)^2\geq0$. Equality occurs at $t=1/2$.
This shows explicitly that eigenvalues can leave the
central line while every finite action identity and
the exact metric-curvature bound remain valid.
It is an example of the proved finite mechanism,
not a counterexample involving the arithmetic zeros.

The completed native result is \eqref{SA13}--\eqref{SA43}:
the original polynomial, every trace Taylor coefficient,
the outgoing relation, the full quotient metric and
the arithmetic rank-two allowance are related by
explicit finite identities. These formulas do not
supply a degree-uniform bound or decay of $\epsilon$.

\begin{thebibliography}{9}\raggedright
\bibitem{CV} Alain Connes and Walter D. van Suijlekom,
\emph{Quadratic Forms, Real Zeros and Echoes of the Spectral Action},
\href{https://arxiv.org/abs/2511.23257v1}{arXiv:2511.23257v1},
28 November 2025. Original author LaTeX
\texttt{Araki-final-oct25.tex}, Section \texttt{sect:sa},
especially \texttt{lem:sa}, its repeated-node proof and
\texttt{prop:der-sa}. This section was read directly for
the present calculation; no PDF was read. Source SHA-256:
\texttt{0357a6123b7b3a2ba2e058458f5267eef1c4ad6abb051af559015b86ae2aa111}.
The source attributes divided differences to Hermite
and the expansional approach to Araki. Those historical
works were not independently read for this finite
polynomial derivation.
\bibitem{Native} Split-Zero programme,
\emph{The native arithmetic secular determinant:
complete pole cancellation and the retained quotient metric},
20 September 2026, \texttt{NATIVE\_SECULAR\_RECEIVER.tex},
NS1--NS37. Its original source measure and cyclic
quotient are retained in \eqref{SA1}--\eqref{SA7};
the finite identities required here are proved again.
\end{thebibliography}
\end{document}
